% -----------------------------------------------------------
% บทที่ 1: บทนำ
% -----------------------------------------------------------

% -----------------------------------------------------------
\section{ความเป็นมาและความสำคัญของโครงงาน}
% -----------------------------------------------------------

หลายคนเคยตัดผมหรือทำสีผมแล้วได้ผลไม่ตรงกับที่คิดไว้ และเมื่อตัดหรือทำสีไปแล้วก็แก้กลับได้ยาก
ถ้ามีระบบที่ช่วยให้เห็นภาพตัวเองในทรงผมและสีผมใหม่ก่อนลงมือจริง ก็จะช่วยให้ตัดสินใจได้ง่ายขึ้น
งานลักษณะนี้ในสาขาคอมพิวเตอร์วิทัศน์ (Computer Vision) เรียกว่าการถ่ายโอนทรงผมและสีผม
(Hair Transfer)~\cite{Zhu_2021,nikolaev2024hairfastganrealisticrobusthair}
ระบบจะรับภาพ 3 ภาพ ได้แก่ ภาพใบหน้าของผู้ใช้ ภาพที่มีทรงผมที่ต้องการ และภาพที่มีสีผมที่ต้องการ
แล้วสร้างภาพใหม่ที่ยังเป็นผู้ใช้คนเดิม สวมเสื้อตัวเดิม และอยู่หน้าฉากหลังเดิม แต่มีทรงผมและสีผมตามภาพอ้างอิง

งานนี้ฟังดูเหมือนเป็นการตัดผมจากภาพหนึ่งไปแปะลงอีกภาพ แต่ในทางปฏิบัติยากกว่านั้นมาก
เพราะเส้นผมมีรายละเอียดเล็ก ๆ จำนวนมาก แสงและเงาบนผมเปลี่ยนไปตามทิศทางของเส้นผม
และขอบผมยังซ้อนทับกับหน้าผาก หู และฉากหลัง
เมื่อเปลี่ยนจากผมยาวเป็นผมสั้น ระบบยังต้องคาดเดาว่าส่วนที่ผมเดิมเคยบังไว้ เช่น หูหรือลำคอ ควรมีลักษณะอย่างไร
ปัญหาอีกข้อหนึ่งคือไม่มีภาพเฉลยให้เทียบ เพราะไม่มีภาพจริงของคนคนนั้นในทรงผมใหม่
การตัดสินว่าภาพใดดีกว่าจึงต้องแยกพิจารณาทีละด้าน เช่น สีตรงหรือไม่ ทรงตรงหรือไม่ และใบหน้ายังเหมือนเดิมหรือไม่
และสุดท้ายยังต้องให้คนที่ดูภาพเป็นผู้ประเมินด้วย

งานวิจัยที่ผ่านมาแก้ปัญหานี้ได้หลายแนวทาง เช่น ใช้โมเดลตระกูล StyleGAN~\cite{karras2020analyzingimprovingimagequality}
ใช้ข้อความสั่งงานผ่านโมเดล CLIP (Contrastive Language-Image Pre-training)~\cite{wei2022hairclipdesignhairtext}
หรือใช้แบบจำลองดิฟฟิวชัน (Diffusion Models)~\cite{rombach2022highresolutionimagesynthesislatent}
แต่ยังไม่มีวิธีใดที่ดีครบทุกด้าน ผู้วิจัยสนใจโมเดล 2 ตัวที่ทำงานต่างกันอย่างชัดเจน
ตัวแรกคือ Barbershop~\cite{Zhu_2021} ซึ่งค้นหารหัสที่ใช้แทนภาพแต่ละภาพด้วยการปรับค่าซ้ำหลายร้อยรอบ
(Optimization-based) จึงได้ภาพที่สมจริง แต่ใช้เวลาหลายนาทีต่อภาพ
ตัวที่สองคือ HairFastGAN~\cite{nikolaev2024hairfastganrealisticrobusthair} ซึ่งใช้ตัวเข้ารหัส (Encoder-based)
คำนวณรหัสออกมาในครั้งเดียว จึงทำงานเสร็จในระดับวินาที
อย่างไรก็ตาม เมื่อผู้วิจัยทดลองใช้ HairFastGAN กับภาพของผู้วิจัยเองในช่วงแรก
พบว่าสีผมมักออกมาเป็นปื้นสีเดียว ไม่ไล่เฉดเหมือนผมจริง และใบหน้าดูซีดกว่าภาพต้นฉบับ

จากข้อสังเกตนี้จึงเกิดคำถามว่า หากนำโมเดลทั้งสองมาปรับจูน (Fine-tune) ด้วยข้อมูลของผู้วิจัยเอง
จุดอ่อนเหล่านี้จะลดลงได้มากเพียงใด และเมื่อปรับแล้วโมเดลใดเหมาะกับการใช้งานจริงมากกว่า
คำถามนี้ตอบจากงานวิจัยเดิมไม่ได้ เพราะแต่ละงานฝึกและวัดผลด้วยข้อมูลและตัวชี้วัดคนละชุด
ผู้วิจัยจึงปรับจูนโมเดลทั้งสองด้วยข้อมูล ฟังก์ชันสูญเสีย และกติกาการเลือกโมเดลชุดเดียวกัน
เพื่อให้การเปรียบเทียบเป็นธรรม

อีกปัญหาหนึ่งที่พบคือ ภาพจากโมเดลโครงข่ายปรปักษ์เชิงกำเนิด (Generative Adversarial Network: GAN) ทั้งสองตัว
มีเส้นผมที่ดูนุ่มและเบลอเกินไป ไม่เห็นเป็นเส้นชัดเหมือนภาพถ่ายจริง
ขณะที่ Stable Diffusion วาดรายละเอียดได้ดีมาก แต่ถ้าให้ Stable Diffusion สร้างทรงผมเองทั้งหมด
จะควบคุมทรงผม สีผม และใบหน้าได้ยาก
ผู้วิจัยจึงแบ่งหน้าที่ให้แต่ละโมเดลทำในส่วนที่ถนัด คือให้ GAN จัดทรงและสีผมให้ถูกต้องก่อน
แล้วใช้ Stable Diffusion แบบแรงต่ำเติมเฉพาะพื้นผิวของเส้นผมและผิวหน้าในภายหลัง
จากนั้นวัดผลทั้งด้วยตัวชี้วัดเชิงตัวเลข 5 ด้าน และด้วยความเห็นของผู้ใช้ที่ไม่รู้ว่าภาพมาจากโมเดลใด
แล้วนำตัวเลือกที่เหมาะสมที่สุดไปพัฒนาเป็นเว็บแอปพลิเคชันให้คนทั่วไปทดลองใช้ได้จริง
เนื่องจากงานวิจัยส่วนใหญ่เผยแพร่เพียงโค้ดต้นฉบับ (Source Code) ซึ่งผู้ที่ไม่ได้เขียนโปรแกรมนำไปใช้ได้ยาก


% -----------------------------------------------------------
\section{วัตถุประสงค์ของโครงงาน}
% -----------------------------------------------------------

\begin{enumerate}
    \item เพื่อศึกษาและเปรียบเทียบประสิทธิภาพของโมเดลการเรียนรู้เชิงลึก 2 โมเดล ได้แก่ HairFastGAN และ Barbershop
    ในการถ่ายโอนทรงผมและสีผม ทั้งก่อนและหลังปรับจูนด้วยข้อมูลและกติกาเดียวกัน
    \item เพื่อวิเคราะห์ข้อดีและข้อจำกัดของแต่ละโมเดล
    ในด้านความถูกต้องของสีผมและทรงผม การรักษาใบหน้าให้เป็นคนเดิม ความสมจริง และเวลาประมวลผล
    รวมถึงผลของการขัดเกลาภาพด้วย Diffusion
    โดยใช้ตัวชี้วัดเชิงตัวเลข 5 ด้านร่วมกับการประเมินโดยผู้ใช้
    \item เพื่อพัฒนาเว็บแอปพลิเคชันต้นแบบที่ให้ผู้ใช้ทดลองเปลี่ยนทรงผมและสีผมได้ง่าย
    โดยใช้โมเดลที่ได้ผลดีที่สุดจากการประเมิน
\end{enumerate}

% -----------------------------------------------------------
\section{ขอบเขตของโครงงาน}
% -----------------------------------------------------------

\begin{enumerate}
    \item ศึกษาและเปรียบเทียบโมเดลหลัก 2 โมเดล ได้แก่ HairFastGAN และ Barbershop
    โดยเริ่มจากน้ำหนักที่ผู้พัฒนาเผยแพร่ไว้ (Pretrained Weights)
    ทั้งสองโมเดลใช้ StyleGAN2 (Style-based Generative Adversarial Network รุ่นที่ 2)~\cite{karras2020analyzingimprovingimagequality}
    ที่ฝึกด้วยชุดภาพใบหน้า FFHQ (Flickr-Faces-HQ) เป็นตัวสร้างภาพ
    \item ใช้ชุดข้อมูลของผู้วิจัย 169 ชุด แต่ละชุดมี 3 ภาพ คือ ภาพใบหน้า ภาพทรงผม และภาพสีผม
    แบ่งเป็นชุดฝึก 150 ชุด ชุดตรวจสอบ 9 ชุด และชุดทดสอบ 10 ชุด
    ภาพทั้งหมดถูกครอปและจัดตำแหน่งใบหน้าตามรูปแบบของ FFHQ ที่ขนาด $1024 \times 1024$ พิกเซล
    \item ปรับจูนเฉพาะส่วนของโมเดลที่ควบคุมสีและรายละเอียดของภาพ
    คือโมดูล Color และ Refinement ของ HairFastGAN
    และชั้นของ StyleGAN2 ที่วาดภาพความละเอียด 64--1024 พิกเซลของ Barbershop
    ทั้งสองโมเดลใช้ฟังก์ชันสูญเสียหลักชุดเดียวกันซึ่งวัด 5 ด้าน
    โดย HairFastGAN มีค่าสูญเสียส่วนเพิ่มที่ใช้กับโมเดลนี้เท่านั้น
    ทดลองอัตราการเรียนรู้ 4 ชุดต่อโมเดล ชุดละ 30 Epoch
    และเลือก Checkpoint ด้วยกติกาเดียวกันบนชุดตรวจสอบ
    \item ขัดเกลาภาพด้วย Stable Diffusion Inpainting ที่ความแรง (Strength) 0.15
    เฉพาะบริเวณผมและใบหน้า แล้วตามด้วยการเก็บรายละเอียดพื้นหลังและการเพิ่มความคม
    โดยไม่ฝึก Diffusion เพิ่ม
    \item เปรียบเทียบ 4 เงื่อนไขหลัก ได้แก่ HairFastGAN B24, HairFastGAN B24 + Diffusion,
    Barbershop-FT A30 และ Barbershop-FT A30 + Diffusion
    (B24 คือน้ำหนักของ Trial B Epoch 24, Barbershop-FT คือ Barbershop ที่ปรับจูนแล้ว
    และ A30 คือน้ำหนักของ Trial A Epoch 30)
    และใช้โมเดลก่อนปรับจูนของทั้งสองเป็นจุดอ้างอิง
    \item ประเมินผลด้วยตัวชี้วัดเชิงตัวเลข 5 ด้าน ได้แก่ ความคลาดเคลื่อนของสีผมในปริภูมิสี CIELAB (Lab),
    ความตรงของทรงผม (Intersection over Union: IoU), ความเหมือนของใบหน้า (ArcFace: Additive Angular Margin Loss),
    การเปลี่ยนแปลงนอกบริเวณผม (Mean Absolute Error: MAE และ Structural Similarity Index: SSIM)
    และรายละเอียดของเส้นผม พร้อมวัดเวลาประมวลผลต่อภาพ
    ร่วมกับการประเมินโดยผู้ใช้ 41 คนแบบไม่บอกชื่อโมเดล
    และการทดลองใช้งานเว็บแอปพลิเคชันจริง 100 ครั้ง
    \item ศึกษาเฉพาะภาพใบหน้าที่หันตรง (Frontal View)
    และประมวลผลการทดลองบน Google Colab ด้วยหน่วยประมวลผลกราฟิก (Graphics Processing Unit: GPU) NVIDIA A100 หน่วยความจำ 40 GB
\end{enumerate}

% -----------------------------------------------------------
\section{ข้อจำกัดของโครงงาน}
% -----------------------------------------------------------
\begin{enumerate}
    \item ชุดตรวจสอบและชุดทดสอบมีขนาดเล็ก (9 และ 10 ชุด)
    จึงใช้ตัวชี้วัดที่ต้องอาศัยภาพจำนวนมาก เช่น FID (Fréchet Inception Distance) ไม่ได้
    และผลของตัวชี้วัดบางตัวยังแกว่งมาก
    \item ระบบต้องใช้ GPU ประสิทธิภาพสูงจึงจะได้เวลาประมวลผลประมาณ 11--14 วินาทีต่อภาพ
    \item ความสามารถในการลบผมเดิมที่บังใบหน้า (Face Occlusion) ขึ้นอยู่กับโมเดลตั้งต้น
    จึงอาจยังเห็นร่องรอยของผมเดิมปนอยู่ในภาพผลลัพธ์
    \item ขั้น Diffusion อาจทำให้รายละเอียดของใบหน้าเปลี่ยนไปเล็กน้อย
    และบริเวณที่ถูกแก้ไขขึ้นกับความแม่นยำของการแบ่งส่วนภาพ
    \item ระบบใช้ได้ดีกับภาพใบหน้าตรงเท่านั้น ผมต้องไม่บังใบหน้า
    ฉากหลังควรเรียบและไม่มีสิ่งใดบัง และภาพต้องถ่ายในที่ที่มีแสงสว่างเพียงพอ
    \item โหมดกล้องบนเว็บแอปพลิเคชันยังไม่เป็นแบบเรียลไทม์เต็มรูปแบบ
    เพราะหลังถ่ายภาพต้องรอประมวลผลประมาณ 15 วินาที ทรงผมใหม่จึงจะปรากฏบนภาพจากกล้อง
\end{enumerate}
